\section{The Chern character and tensor products}
\label{section-chern-classes-tensor}

\noindent
Let $(S, \delta)$ be as in Situation \ref{situation-setup}.
Let $X$ be locally of finite type over $S$.
We define the {\it Chern character} of a finite locally free
$\mathcal{O}_X$-module to be the formal expression
$$
ch({\mathcal E}) = \sum\nolimits_{i=1}^r e^{x_i}
$$
if the $x_i$ are the Chern roots of ${\mathcal E}$. Writing this
as a polynomial in the Chern classes we obtain
\begin{align*}
ch(\mathcal{E})
& =
r(\mathcal{E})
+
c_1(\mathcal{E}) +
\frac{1}{2}(c_1(\mathcal{E})^2 - 2c_2(\mathcal{E}))
+
\frac{1}{6}(c_1(\mathcal{E})^3 - 3c_1(\mathcal{E})c_2(\mathcal{E}) + 3c_3(\mathcal{E})) \\
& \quad\quad +
\frac{1}{24}(c_1(\mathcal{E})^4 - 4c_1(\mathcal{E})^2c_2(\mathcal{E}) + 4c_1(\mathcal{E})c_3(\mathcal{E}) + 2c_2(\mathcal{E})^2 - 4c_4(\mathcal{E}))
+
\ldots \\
& =
\sum\nolimits_{p = 0, 1, 2, \ldots} \frac{P_p(\mathcal{E})}{p!}
\end{align*}
with $P_p$ polynomials in the Chern classes as in
Example \ref{example-power-sum}. The degree $p$ component of
the above is
$$
ch_p(\mathcal{E}) = \frac{P_p(\mathcal{E})}{p!} \in A^p(X) \otimes \mathbf{Q}
$$
What does it mean that the coefficients are rational numbers?
Well this simply means that we think of
$ch_p(\mathcal{E})$ as an element of $A^p(X) \otimes \mathbf{Q}$.

\begin{remark}
\label{remark-extend-chern-character-to-finite-locally-free}
In the discussion above we have defined the components of the Chern character
$ch_p(\mathcal{E}) \in A^p(X) \otimes \mathbf{Q}$
of $\mathcal{E}$ even if the rank of $\mathcal{E}$
is not constant. See Remarks \ref{remark-extend-to-finite-locally-free} and
\ref{remark-equalities-nonconstant-rank}. Thus the full Chern character
of $\mathcal{E}$ is
an element of $\prod_{p \geq 0} (A^p(X) \otimes \mathbf{Q})$. If $X$
is quasi-compact and $\dim(X) < \infty$ (usual dimension), then one can show
using Lemma \ref{lemma-vanish-above-dimension} and the splitting principle
that $ch(\mathcal{E}) \in A^*(X) \otimes \mathbf{Q}$.
\end{remark}

\begin{lemma}
\label{lemma-chern-character-additive}
Let $(S, \delta)$ be as in Situation \ref{situation-setup}. Let $X$ be locally
of finite type over $S$. Let
$
0 \to \mathcal{E}_1 \to \mathcal{E} \to \mathcal{E}_2 \to 0
$
be a short exact sequence of finite locally free $\mathcal{O}_X$-modules.
Then we have the equality
$$
ch(\mathcal{E}) = ch(\mathcal{E}_1) + ch(\mathcal{E}_2)
$$
More precisely, we have
$P_p(\mathcal{E}) = P_p(\mathcal{E}_1) + P_p(\mathcal{E}_2)$
in $A^p(X)$ where $P_p$ is as in Example \ref{example-power-sum}.
\end{lemma}

\begin{proof}
It suffices to prove the more precise statement. By
Section \ref{section-splitting-principle}
this follows because if $x_{1, i}$, $i = 1, \ldots, r_1$
and $x_{2, i}$, $i = 1, \ldots, r_2$ are the
Chern roots of $\mathcal{E}_1$ and $\mathcal{E}_2$, then
$x_{1, 1}, \ldots, x_{1, r_1}, x_{2, 1}, \ldots, x_{2, r_2}$
are the Chern roots of $\mathcal{E}$. Hence we get the result
from our choice of $P_p$ in Example \ref{example-power-sum}.
\end{proof}

\begin{lemma}
\label{lemma-chern-character-multiplicative}
Let $(S, \delta)$ be as in Situation \ref{situation-setup}. Let $X$ be locally
of finite type over $S$. Let $\mathcal{E}_1$ and $\mathcal{E}_2$
be finite locally free $\mathcal{O}_X$-modules.
Then we have the equality
$$
ch(\mathcal{E}_1 \otimes_{\mathcal{O}_X} \mathcal{E}_2) =
ch(\mathcal{E}_1) ch(\mathcal{E}_2)
$$
More precisely, we have
$$
P_p(\mathcal{E}_1 \otimes_{\mathcal{O}_X} \mathcal{E}_2) =
\sum\nolimits_{p_1 + p_2 = p}
{p \choose p_1} P_{p_1}(\mathcal{E}_1) P_{p_2}(\mathcal{E}_2)
$$
in $A^p(X)$ where $P_p$ is as in Example \ref{example-power-sum}.
\end{lemma}

\begin{proof}
It suffices to prove the more precise statement. By
Section \ref{section-splitting-principle}
this follows because if $x_{1, i}$, $i = 1, \ldots, r_1$
and $x_{2, i}$, $i = 1, \ldots, r_2$ are the
Chern roots of $\mathcal{E}_1$ and $\mathcal{E}_2$, then
$x_{1, i} + x_{2, j}$, $1 \leq i \leq r_1$, $1 \leq j \leq r_2$
are the Chern roots of $\mathcal{E}_1 \otimes \mathcal{E}_2$.
Hence we get the result from the binomial formula for
$(x_{1, i} + x_{2, j})^p$ and the
shape of our polynomials $P_p$ in Example \ref{example-power-sum}.
\end{proof}

\begin{lemma}
\label{lemma-chern-character-dual}
In Situation \ref{situation-setup} let $X$ be locally of finite type over $S$.
Let $\mathcal{E}$ be a finite locally free $\mathcal{O}_X$-module
with dual $\mathcal{E}^\vee$. Then
$ch_i(\mathcal{E}^\vee) = (-1)^i ch_i(\mathcal{E})$ in
$A^i(X) \otimes \mathbf{Q}$.
\end{lemma}

\begin{proof}
Follows from the corresponding result for Chern classes
(Lemma \ref{lemma-chern-classes-dual}).
\end{proof}











